\chapter{Continuous measurement: the Zeno effect and quantum trajectories}
\label{ch:continuous-measurement}

The master equation of \cref{ch:open-systems} describes the average
state of an open system: average over the environment, average over
all the things the bath might have recorded. But the input--output
theory of \cref{ch:input-output} showed that the bath's output
field can be \emph{measured} --- a detector downstream turns it into
a record. What happens to the system state when we actually look?
Two phenomena emerge that the averaged master equation hides. First,
measurement that is frequent enough freezes the dynamics: the
\emph{quantum Zeno effect}. Second, a single experimental run does
not show the smooth exponential decay of the master equation; it
shows abrupt \emph{quantum jumps} at random times, and the master
equation is only what one recovers after averaging many such runs.
This closing chapter of Part~II treats both. We derive the Zeno
freezing from the short-time behaviour of the survival probability,
then unravel the Lindblad equation into individual
\emph{quantum trajectories} --- the stochastic conditioned
evolution of a continuously monitored system --- and show that
averaging the trajectories returns the master equation.

% --------------------------------------------------------------
\section{The quantum Zeno effect}
\label{sec:zeno}

\begin{phenomenon}
A watched quantum system does not evolve. If a system that would
otherwise leave its initial state is measured repeatedly --- each
measurement asking ``are you still in the initial state?'' --- then
in the limit of infinitely frequent measurement the system is
frozen in place, never decaying. The effect is not mystical: it
follows from the fact that the survival probability of a quantum
state decays \emph{quadratically}, not linearly, at short times, so
that resetting the clock often enough never lets the quadratic
decay accumulate.
\end{phenomenon}

\begin{statement}[\textbf{Quantum Zeno effect}~\cite{misra1977}]
\label{stmt:zeno}
Let a system start in $\ket{\psi_0}$ and evolve under $\op H$. The
short-time survival probability is
\begin{equation}
\label{eq:survival-short}
P(\tau) \;=\; \bigl|\bra{\psi_0}e^{-i\op H\tau/\hbar}\ket{\psi_0}\bigr|^2
\;=\; 1 - \Bigl(\frac{\tau}{\tau_Z}\Bigr)^2 + \mathcal O(\tau^4),
\qquad
\frac{1}{\tau_Z^2} \;=\; \frac{(\Delta H)^2}{\hbar^2},
\end{equation}
where $(\Delta H)^2=\bra{\psi_0}\op H^2\ket{\psi_0}-\bra{\psi_0}\op H\ket{\psi_0}^2$
is the energy variance in the initial state. If the state is
projectively measured $N$ times during a total interval $t$, at
spacing $\tau=t/N$, the probability of surviving all $N$
measurements is
\begin{equation}
\label{eq:zeno-survival}
P_N(t) \;=\; \Bigl[1 - \Bigl(\frac{t}{N\tau_Z}\Bigr)^2\Bigr]^N
\;\xrightarrow[N\to\infty]{}\; 1.
\end{equation}
Infinitely frequent measurement freezes the system in
$\ket{\psi_0}$.
\end{statement}

\paragraph{Derivation.}
Expand the propagator to second order,
$e^{-i\op H\tau/\hbar}\ket{\psi_0}=\ket{\psi_0}-\tfrac{i\tau}{\hbar}\op H\ket{\psi_0}-\tfrac{\tau^2}{2\hbar^2}\op H^2\ket{\psi_0}+\cdots$.
The amplitude is
$\bra{\psi_0}e^{-i\op H\tau/\hbar}\ket{\psi_0}=1-\tfrac{i\tau}{\hbar}\langle\op H\rangle-\tfrac{\tau^2}{2\hbar^2}\langle\op H^2\rangle+\cdots$,
and its modulus squared is
$P(\tau)=1-\tfrac{\tau^2}{\hbar^2}(\langle\op H^2\rangle-\langle\op H\rangle^2)+\mathcal O(\tau^4)$,
the linear term cancelling --- this is the crucial point. With
$N$ measurements each of duration $\tau=t/N$, the survival
probability multiplies, $P_N(t)=P(t/N)^N=[1-(t/N\tau_Z)^2]^N$.
Taking the logarithm, $\ln P_N = N\ln[1-(t/N\tau_Z)^2]\approx -t^2/(N\tau_Z^2)\to0$
as $N\to\infty$, so $P_N\to1$. The quadratic short-time decay is
what makes the product go to one; a genuinely linear (exponential)
decay would give $P_N=(1-t/N\tau)^N\to e^{-t/\tau}$, no freezing.

\paragraph{Continuous measurement and the anti-Zeno effect.}
Replacing discrete projective measurements by continuous monitoring
(strong coupling to a meter or to a fast-decaying readout mode)
produces the same freezing: a continuously watched transition is
suppressed. The first clean demonstration was the trapped-ion
experiment of Itano~et~al.~\cite{itano1990}, in which frequent
measurement pulses inhibited a driven hyperfine transition. The
effect has a partner --- the \emph{anti-Zeno} effect --- in which
measurement at an intermediate rate \emph{accelerates} decay,
because the measurement bandwidth then overlaps the part of the
environment spectrum that drives the transition. Which one occurs
depends on the bath spectral density of \cref{sec:noise-psd}
sampled by the measurement.

\paragraph{Transition.}
The Zeno effect is the strong-measurement limit. To describe
measurement of \emph{any} strength --- and to see what a single
experimental run actually looks like --- we unravel the master
equation into trajectories conditioned on the measurement record.

% --------------------------------------------------------------
\section{Unravelling the master equation}
\label{sec:unravelling}

\begin{phenomenon}
A fluorescing atom monitored by a perfect photodetector is either
emitting or not: the record is a train of discrete clicks at random
times, and between clicks the atom's state evolves smoothly. The
excited-state population of a \emph{single} atom does not decay
smoothly --- it sits, then jumps to the ground state at the instant
a photon is detected. The smooth exponential decay of the master
equation is the \emph{average} over many such jump records. An
\emph{unravelling} is a stochastic equation for the conditioned
state of one run, constructed so that its ensemble average
reproduces the Lindblad equation.
\end{phenomenon}

\begin{statement}[\textbf{Quantum-jump (Monte Carlo wavefunction) unravelling}~\cite{dalibard1992,carmichael1993}]
\label{stmt:quantum-jumps}
Consider a system with master equation
$\dot{\op\rho}=-\tfrac{i}{\hbar}\comm{\op H}{\op\rho}+\sum_k\gamma_k\mathcal D[\op L_k]\op\rho$
whose output channels are perfectly photon-counted. In a time step
$\dd t$, the conditioned pure state $\ket{\psi}$ either jumps or
does not. A jump in channel $k$ occurs with probability
\begin{equation}
\label{eq:jump-prob}
p_k \;=\; \gamma_k\,\dd t\,\bra{\psi}\op L_k^\dagger\op L_k\ket{\psi},
\end{equation}
upon which the state is reset,
$\ket{\psi}\to\op L_k\ket{\psi}/\lVert\op L_k\ket{\psi}\rVert$ (a
detector click). With probability $1-\sum_k p_k$ no jump occurs and
the state evolves under the non-Hermitian effective Hamiltonian
\begin{equation}
\label{eq:effective-H}
\op H_{\rm eff} \;=\; \op H - \frac{i\hbar}{2}\sum_k\gamma_k\,\op L_k^\dagger\op L_k,
\end{equation}
followed by renormalisation. Averaging the conditioned state
$\ket{\psi}\!\bra{\psi}$ over the jump/no-jump branches reproduces
the Lindblad master equation exactly.
\end{statement}

\paragraph{Derivation from the Kraus operators.}
The unravelling is read directly off the infinitesimal Kraus
operators of \cref{stmt:kraus}. The jump operators
$\op K_k=\sqrt{\gamma_k\dd t}\,\op L_k$ implement the
detector-click update with probability
$\lVert\op K_k\ket\psi\rVert^2=\gamma_k\dd t\bra\psi\op L_k^\dagger\op L_k\ket\psi=p_k$,
exactly \cref{eq:jump-prob}; the no-jump operator
$\op K_0=\id-(\tfrac{i}{\hbar}\op H+\tfrac12\sum_k\gamma_k\op L_k^\dagger\op L_k)\dd t
=\id-\tfrac{i}{\hbar}\op H_{\rm eff}\dd t$ implements the smooth
non-Hermitian evolution with probability $1-\sum_k p_k$. The
ensemble average over the two branches is, to first order in
$\dd t$,
\begin{equation}
\label{eq:trajectory-average}
\overline{\op\rho(t+\dd t)}
= \op K_0\op\rho\op K_0^\dagger + \sum_k\op K_k\op\rho\op K_k^\dagger
= \op\rho + \dd t\Bigl[-\tfrac{i}{\hbar}\comm{\op H}{\op\rho}+\sum_k\gamma_k\mathcal D[\op L_k]\op\rho\Bigr],
\end{equation}
which is the Lindblad equation. The master equation is precisely
the average over all measurement records; the trajectory is one
sample of that average.

\paragraph{The non-Hermitian no-jump evolution is informative.}
The decay term in $\op H_{\rm eff}$ is not merely bookkeeping: even
when no photon is detected, the state changes, because the
\emph{absence} of a click is itself information. For a qubit
monitored on $\op L=\hat\sigma_-$, the longer one waits without
seeing a photon, the more probable it is that the qubit was already
in the ground state, and the no-jump evolution steadily increases
the ground-state weight. This measurement back-action with no
detector event is a genuinely quantum feature with no classical
analogue.

\begin{statement}[\textbf{Quantum-state-diffusion (homodyne) unravelling}~\cite{wiseman2009}]
\label{stmt:qsd}
If the output field is instead measured by homodyne detection ---
mixing it with a strong local oscillator and recording a
quadrature --- the conditioned state evolves not by discrete jumps
but by a continuous diffusion. For a single channel $\op L$ the
stochastic master equation is
\begin{equation}
\label{eq:sme}
\dd\op\rho_c
= -\frac{i}{\hbar}\comm{\op H}{\op\rho_c}\dd t
+ \gamma\,\mathcal D[\op L]\op\rho_c\,\dd t
+ \sqrt{\gamma\,\eta}\,\mathcal H[\op L]\op\rho_c\,\dd W,
\end{equation}
where $\dd W$ is a Wiener increment (Gaussian white noise with
$\overline{\dd W}=0$, $\dd W^2=\dd t$), $\eta$ is the detector
efficiency, and the measurement superoperator is
$\mathcal H[\op L]\op\rho = \op L\op\rho+\op\rho\op L^\dagger
-\langle\op L+\op L^\dagger\rangle\op\rho$. Averaging over the noise
($\overline{\dd W}=0$) again returns the Lindblad equation.
\end{statement}

\paragraph{Two unravellings, one master equation.}
Quantum jumps and quantum state diffusion are two different
\emph{measurement schemes} applied to the \emph{same} open system:
photon counting gives a jumpy record, homodyne gives a diffusive
one. They produce genuinely different conditioned trajectories ---
the qubit's path through state space depends on how you look --- yet
both average to the same Lindblad equation. The unravelling is not
unique; it is fixed by the choice of detector, exactly as the Kraus
decomposition of \cref{stmt:kraus} is non-unique.

% --------------------------------------------------------------
\section{Quantum jumps in the laboratory}
\label{sec:jumps-lab}

\begin{example}[\textbf{Telegraph record of a monitored qubit}]
\label{ex:telegraph}
A qubit dispersively coupled to a continuously monitored readout
cavity (\cref{ch:input-output}) realises a continuous measurement
of $\hat\sigma_z$. The homodyne record then tracks the qubit state
in real time, and quantum jumps between $\ket0$ and $\ket1$ appear
as a random telegraph signal. The same machinery, applied to the
$\ket{e}\to\ket{f}$ monitoring of a three-level transmon, underlies
the ``catch and reverse'' experiments that detect a quantum jump
\emph{as it begins} and reverse it mid-flight --- a direct
laboratory use of the no-jump back-action of
\cref{eq:effective-H}. The photon-counting unravelling
(\cref{stmt:quantum-jumps}) is also the natural language for the
single-microwave-photon detectors of \cref{ch:microwave-absorption}:
each engineered Lindblad jump $\op L=\sqrt{\kappa_{\rm nl}}\,\hat b\,\hat\sigma^\dagger$
is precisely a detector click registering the arrival of a signal
photon.
\end{example}

\paragraph{Why trajectories matter for sensing.}
For the dark-matter sensing of Part~VI, the distinction between the
averaged master equation and the conditioned trajectory is the
distinction between the \emph{ensemble-averaged} dark-count rate and
the \emph{single-shot} click record actually used to set a limit.
The detector is read out one trajectory at a time; the master
equation supplies the average rates that calibrate it, while the
trajectory picture supplies the statistics of individual clicks.
The continuous-measurement formalism is therefore not a
mathematical luxury --- it is the language in which a
single-quantum detector is actually operated.

% --------------------------------------------------------------
This chapter completed Part~II by asking what happens when the
environment is not discarded but measured. Frequent measurement
freezes the dynamics --- the quantum Zeno effect --- because the
survival probability decays quadratically at short times.
Measurement of arbitrary strength is described by unravelling the
master equation into quantum trajectories: photon counting gives
the quantum-jump unravelling, with discrete clicks and informative
no-jump evolution, while homodyne detection gives the diffusive
stochastic master equation. Both average back to the Lindblad
equation of \cref{ch:open-systems}.

This closes the open-systems toolkit. Part~III builds on it to treat
quantum optics proper --- the quantisation of the electromagnetic
field, the Jaynes--Cummings model, and the dispersive regime ---
where the cavity, the input--output relations, and the
continuous-measurement picture assembled here become the working
description of light--matter interaction and qubit readout.

\clearpage
\begin{chaptersummary}

\paragraph{Quantum Zeno effect.}
The survival probability decays quadratically at short times,
$P(\tau)=1-(\tau/\tau_Z)^2$ with $1/\tau_Z^2=(\Delta H)^2/\hbar^2$.
With $N$ measurements in time $t$,
$P_N(t)=[1-(t/N\tau_Z)^2]^N\to1$ as $N\to\infty$: frequent
measurement freezes the system. Intermediate measurement rates can
instead accelerate decay (anti-Zeno), depending on the bath
spectrum.

\paragraph{Quantum-jump unravelling.}
A photon-counted system jumps with probability
$p_k=\gamma_k\dd t\langle\op L_k^\dagger\op L_k\rangle$ (state reset
$\ket\psi\to\op L_k\ket\psi/\lVert\cdot\rVert$); otherwise it evolves
under $\op H_{\rm eff}=\op H-\tfrac{i\hbar}{2}\sum_k\gamma_k\op L_k^\dagger\op L_k$.
The no-jump evolution is informative (absence of a click updates the
state). Averaging the branches gives the Lindblad equation.

\paragraph{Quantum-state diffusion.}
Homodyne detection gives the stochastic master equation
$\dd\op\rho_c=-\tfrac{i}{\hbar}\comm{\op H}{\op\rho_c}\dd t+\gamma\mathcal D[\op L]\op\rho_c\dd t+\sqrt{\gamma\eta}\,\mathcal H[\op L]\op\rho_c\,\dd W$
with Wiener noise $\dd W$. Different detectors give different
unravellings of the same master equation; the choice is fixed by
how the output field is measured.

\paragraph{Laboratory picture.}
A continuously monitored qubit shows a random-telegraph jump record;
catch-and-reverse exploits the no-jump back-action; each engineered
Lindblad jump in a single-photon detector is a detector click. The
master equation supplies average rates; the trajectory supplies the
single-shot statistics.

\end{chaptersummary}

% --------------------------------------------------------------
\section*{Exercises}
\addcontentsline{toc}{section}{Exercises}

\begin{exercise}[Quadratic short-time decay]
Derive $P(\tau)=1-(\Delta H)^2\tau^2/\hbar^2+\mathcal O(\tau^4)$ from
the second-order expansion of the propagator, and show explicitly
that the linear-in-$\tau$ term cancels.
\begin{solution}
$\bra{\psi_0}e^{-i\op H\tau/\hbar}\ket{\psi_0}
=1-\tfrac{i\tau}{\hbar}\langle\op H\rangle-\tfrac{\tau^2}{2\hbar^2}\langle\op H^2\rangle+\cdots$.
Its modulus squared:
$P=|1-\tfrac{i\tau}\hbar\langle\op H\rangle-\tfrac{\tau^2}{2\hbar^2}\langle\op H^2\rangle|^2
=1-\tfrac{\tau^2}{\hbar^2}\langle\op H^2\rangle+\tfrac{\tau^2}{\hbar^2}\langle\op H\rangle^2+\mathcal O(\tau^4)$
(the $\mathcal O(\tau)$ pieces are imaginary and cancel in the
modulus). Hence $P=1-\tfrac{\tau^2}{\hbar^2}(\langle\op H^2\rangle-\langle\op H\rangle^2)$,
quadratic with no linear term.
\end{solution}
\end{exercise}

\begin{exercise}[Survival in the Zeno limit]
Show that $P_N(t)=[1-(t/N\tau_Z)^2]^N\to1$ as $N\to\infty$, and
contrast with the result for genuinely exponential decay.
\begin{solution}
$\ln P_N=N\ln[1-(t/N\tau_Z)^2]\approx -N\cdot(t/N\tau_Z)^2=-t^2/(N\tau_Z^2)\to0$,
so $P_N\to1$. For linear decay $P(\tau)=1-\tau/\tau$:
$P_N=(1-t/N\tau)^N\to e^{-t/\tau}$, a finite survival probability,
no freezing. The quadratic short-time law is essential to the Zeno
effect.
\end{solution}
\end{exercise}

\begin{exercise}[Zeno time of a driven qubit]
A qubit driven on resonance has $\op H=\tfrac{\hbar\Omega}{2}\hat\sigma_x$
and starts in $\ket0$. Compute $\tau_Z$ and the measurement spacing
needed to suppress the Rabi transition.
\begin{solution}
$(\Delta H)^2=\langle\op H^2\rangle-\langle\op H\rangle^2$ in
$\ket0$: $\op H^2=(\hbar\Omega/2)^2\id$ so $\langle\op H^2\rangle=(\hbar\Omega/2)^2$,
and $\langle\op H\rangle=\tfrac{\hbar\Omega}{2}\bra0\hat\sigma_x\ket0=0$.
Hence $(\Delta H)^2=(\hbar\Omega/2)^2$ and $\tau_Z=2/\Omega$.
Measurement spacing $\tau\ll\tau_Z=2/\Omega$, i.e.\ much faster than
the Rabi period, freezes the qubit in $\ket0$.
\end{solution}
\end{exercise}

\begin{exercise}[Jump probability and normalisation]
Show that the total jump probability $\sum_k p_k$ plus the no-jump
weight $\lVert\op K_0\ket\psi\rVert^2$ equals $1$ to first order in
$\dd t$, so the unravelling conserves probability.
\begin{solution}
$\lVert\op K_0\ket\psi\rVert^2=\bra\psi\op K_0^\dagger\op K_0\ket\psi
=\bra\psi(\id-\sum_k\gamma_k\op L_k^\dagger\op L_k\dd t)\ket\psi
=1-\sum_k\gamma_k\dd t\langle\op L_k^\dagger\op L_k\rangle=1-\sum_k p_k$
to first order (the $\op H$ part is anti-Hermitian in $\op K_0$ and
drops from $\op K_0^\dagger\op K_0$ at $\mathcal O(\dd t)$). Adding
$\sum_k p_k$ gives $1$.
\end{solution}
\end{exercise}

\begin{exercise}[Averaging reproduces Lindblad]
Verify \cref{eq:trajectory-average}: that
$\op K_0\op\rho\op K_0^\dagger+\sum_k\op K_k\op\rho\op K_k^\dagger$
equals $\op\rho+\dd t\,\mathcal L\op\rho$ to first order in $\dd t$.
\begin{solution}
$\op K_k\op\rho\op K_k^\dagger=\gamma_k\dd t\,\op L_k\op\rho\op L_k^\dagger$.
$\op K_0\op\rho\op K_0^\dagger
=(\id-\tfrac{i}{\hbar}\op H_{\rm eff}\dd t)\op\rho(\id+\tfrac{i}{\hbar}\op H_{\rm eff}^\dagger\dd t)
=\op\rho-\tfrac{i}{\hbar}\dd t(\op H_{\rm eff}\op\rho-\op\rho\op H_{\rm eff}^\dagger)$.
With $\op H_{\rm eff}=\op H-\tfrac{i\hbar}2\sum_k\gamma_k\op L_k^\dagger\op L_k$,
$\op H_{\rm eff}\op\rho-\op\rho\op H_{\rm eff}^\dagger
=\comm{\op H}{\op\rho}-\tfrac{i\hbar}2\sum_k\gamma_k\{\op L_k^\dagger\op L_k,\op\rho\}$.
Substituting and adding the jump terms gives
$\op\rho+\dd t[-\tfrac{i}{\hbar}\comm{\op H}{\op\rho}+\sum_k\gamma_k(\op L_k\op\rho\op L_k^\dagger-\tfrac12\{\op L_k^\dagger\op L_k,\op\rho\})]$,
the Lindblad equation.
\end{solution}
\end{exercise}

\begin{exercise}[No-jump evolution of a decaying qubit]
For $\op L=\hat\sigma_-$ at rate $\gamma$ and no detected photon up
to time $t$, find the conditioned state given
$\ket{\psi(0)}=\alpha\ket0+\beta\ket1$. Show the ground-state weight
grows.
\begin{solution}
$\op H_{\rm eff}=-\tfrac{i\hbar\gamma}2\hat\sigma_+\hat\sigma_-
=-\tfrac{i\hbar\gamma}2\ket1\!\bra1$ (taking $\op H=0$). Then
$e^{-i\op H_{\rm eff}t/\hbar}\ket{\psi(0)}=\alpha\ket0+\beta e^{-\gamma t/2}\ket1$.
After renormalising,
$\ket{\psi(t)}=(\alpha\ket0+\beta e^{-\gamma t/2}\ket1)/\sqrt{|\alpha|^2+|\beta|^2e^{-\gamma t}}$.
The ground-state weight $|\alpha|^2/(|\alpha|^2+|\beta|^2e^{-\gamma t})\to1$:
seeing no photon makes it ever more likely the qubit was already in
$\ket0$ --- back-action with no detector event.
\end{solution}
\end{exercise}

\begin{exercise}[Jump resets the qubit]
For $\op L=\hat\sigma_-$, show that a detected photon resets any
state to $\ket0$, and explain why the post-jump state is
independent of the pre-jump amplitudes.
\begin{solution}
$\op L\ket\psi=\hat\sigma_-(\alpha\ket0+\beta\ket1)=\beta\ket0$;
normalised, $\op L\ket\psi/\lVert\op L\ket\psi\rVert=\ket0$. The
photon carries away the excitation, so detecting it projects the
qubit to the ground state regardless of $\alpha,\beta$ (only the
\emph{probability} of the jump, $\propto|\beta|^2$, depends on the
amplitudes).
\end{solution}
\end{exercise}

\begin{exercise}[Wiener increment algebra]
Using $\dd W^2=\dd t$ and $\overline{\dd W}=0$, show that averaging
the stochastic master equation \cref{eq:sme} over the noise yields
the Lindblad equation, and explain the role of $\eta$.
\begin{solution}
Taking the ensemble average, the stochastic term
$\sqrt{\gamma\eta}\,\mathcal H[\op L]\op\rho_c\,\dd W$ averages to
zero since $\overline{\dd W}=0$ and $\op\rho_c$ is independent of
the future increment. The deterministic terms are exactly the
Lindblad generator, so $\overline{\dd\op\rho_c}=\mathcal L\overline{\op\rho_c}\,\dd t$.
The efficiency $\eta\in[0,1]$ scales only the measurement
(stochastic) term: at $\eta=0$ (no detection) the conditioned state
equals the unconditioned master-equation state; at $\eta=1$ (ideal
detection) the measurement back-action is maximal. The averaged
equation is independent of $\eta$.
\end{solution}
\end{exercise}

\begin{exercise}[Jumps versus diffusion]
Explain why photon counting and homodyne detection of the same
emission channel produce qualitatively different single-run records,
even though both average to the same master equation.
\begin{solution}
Photon counting measures the number operator-like channel
$\op L^\dagger\op L$: the record is a discrete click train, and the
state makes abrupt jumps. Homodyne measures a field quadrature
$\op L+\op L^\dagger$: the record is a continuous noisy current, and
the state diffuses smoothly. The two correspond to different
positive-operator-valued measures on the output field, i.e.\
different unravellings; only the noise-averaged dynamics, which
discards the record, is common to both --- the Lindblad equation.
\end{solution}
\end{exercise}

\begin{exercise}[Anti-Zeno qualitatively]
Explain, using the noise-spectral-density picture of
\cref{sec:noise-psd}, why measurement at an intermediate rate can
\emph{accelerate} decay rather than freeze it.
\begin{solution}
A measurement of rate $1/\tau$ broadens the effective level into a
band of width $\sim1/\tau$. Decay proceeds at a rate set by the
overlap of this measurement band with the environment spectral
density $S(\omega)$ at the transition. Very fast measurement
($1/\tau$ huge) broadens the band far beyond where $S$ has weight,
suppressing decay (Zeno). At an intermediate rate the band can be
tuned to \emph{maximise} overlap with a peak of $S(\omega)$ away
from the bare transition, enhancing the decay rate --- the anti-Zeno
effect. Which dominates depends on the shape of $S(\omega)$.
\end{solution}
\end{exercise}
